\chapter{Why Long-Range Learning Fails and How Models Address It}

Message passing gives source information a route to a target node, but reachability alone does not make that information useful.
A model must preserve task-relevant distinctions during forward propagation, carry enough relevant information through its communication channels, and learn suitable transformations through optimisation.
These requirements motivate three common explanations of long-range failure: over-smoothing, over-squashing, and vanishing gradients \cite{li2018deeper,alon2021bottleneck,glorot2010difficulty}.

The three explanations concern different internal traces.
Over-smoothing concerns contraction of node representations that removes task-relevant distinctions.
Over-squashing concerns inadequate usable influence from relevant source information even though that information is structurally reachable.
Vanishing gradients concern weak backward signals in the parameters that must learn the required computation.
An accuracy decline as dependency length $d$ grows is compatible with all three mechanisms, but it identifies none of them.

Standard message-passing architectures provide useful comparisons because they vary the local operation while retaining one-hop propagation.
GCN uses normalised neighbourhood averaging, whereas GraphSAGE learns inductive neighbourhood aggregation \cite{kipf2017semi,hamilton2017graphsage}.
GAT learns local attention weights, and GIN uses expressive sum aggregation \cite{velickovic2018gat,xu2019gin}.
SGC instead separates repeated propagation from intermediate non-linear transformations \cite{wu2019sgc}.
These differences can change how information survives each hop, but a local architecture still needs sufficient depth for a source node at graph distance $d$ to affect the target node.
Architectures and interventions appear below only when their design addresses a predicted trace or a separate limitation such as strict under-reaching.

This chapter separates mechanisms from symptoms, diagnostics, and interventions.
Sections 2--4 define each mechanism, state its predicted internal trace, group the methods that target it, and delimit what a successful intervention would establish.
Section 5 explains how the mechanisms interact.
Section 6 turns the resulting distinctions into a bounded diagnostic taxonomy for Chapter 6.

\section{Evidence for a Failure Explanation}

A failure explanation must connect an observed performance symptom to a mechanism inside the learned system.
For a dependency of length $d$, the source information must first lie inside the target node's receptive field.
If the model has too few local message-passing layers, the source has exactly zero influence on the target representation.
This case is strict under-reaching, not over-smoothing, over-squashing, or an optimisation failure.

Once the source is reachable, a diagnosis requires a mechanism-specific trace.
Over-smoothing predicts contraction among representations that must remain distinguishable for the task.
Over-squashing predicts inadequate source-to-target influence because topology, dynamics, finite width, or task-relative mixing limits the available communication.
Vanishing gradients predict that the backward signal becomes too small in the layers or parameters that must learn the dependency.
The trace must be measured in the same regime in which the performance symptom appears.

This evidence standard separates four levels of analysis.
An accuracy curve over $d$ is a performance symptom.
A proposed account such as representation contraction is a candidate mechanism.
A measurement such as projection distance, source-to-target sensitivity, or an early-layer gradient norm is a diagnostic.
A residual connection, rewiring rule, or global-attention layer is an intervention.
Conflating these levels turns a method's motivation into a causal conclusion that the experiment has not established.

An intervention can provide stronger evidence when it targets a predicted trace.
The comparison should hold the task, data, effective dependency, model capacity, and training protocol as constant as possible.
The intervention should change the trace in the predicted direction and improve the task outcome in the same runs.
Even this joint result supports a causal account only under the assumptions of the comparison because most interventions change more than one property of the system.

The diagnostic must also be task-aligned.
A graph-level structural proxy may identify a narrow cut without showing that task-relevant information crosses it.
A representation statistic may show contraction without showing that the contracted direction carries the label.
A parameter-gradient norm may be nonzero even when the update direction does not help learn the intended dependency.
The useful question is therefore not whether a familiar pathology can occur, but whether its predicted trace appears in the failing computation and concerns information needed by the target function.

\section{Over-Smoothing}

\subsection{Definition}

Graph propagation often contracts node-varying components of a representation towards an operator-dependent invariant or low-dimensional subspace.
The graph-smoothing interpretation of GCN makes this effect explicit, while later analyses give sufficient spectral and weight-norm conditions for exponential convergence towards an invariant subspace \cite{li2018deeper,oono2020exponentially}.
Over-smoothing occurs when this contraction removes task-relevant distinctions before the model completes the required computation.
The limiting representation need not be a single constant vector, so over-smoothing is broader than literal equality among all node representations.

The mechanism is task-sensitive.
Finite propagation can denoise features or improve a prediction before further smoothing becomes harmful, as shown in model-specific regression and classification analyses \cite{keriven2022not}.
The same amount of contraction can be harmless for a target that depends on a low-frequency graph signal and destructive for a target that requires a distinction between particular nodes.
Depth is therefore not itself evidence of over-smoothing, and an accuracy decrease with depth is only a symptom.

The mechanism is also measurement-sensitive.
Pairwise Euclidean distances depend on representation scale, while cosine similarity ignores changes in norm.
Distance to a theoretically predicted invariant subspace tests a more specific account than a generic variance statistic.
A task-conditioned probe asks whether the remaining representation still contains the distinction required for prediction.
Each diagnostic supports a different statement about contraction.

\subsection{Predicted Internal Trace}

Over-smoothing predicts progressive convergence in the representations of nodes or examples that the target function requires the model to distinguish.
The trace should strengthen with propagation depth or dependency length under otherwise comparable conditions.
It should appear before the task readout, because the proposed mechanism concerns information lost during message passing rather than only a poorly fitted decision boundary.

Complete collapse and label mixing are not the same trace.
Complete collapse maps the measured representations to one value or to the invariant subspace predicted by the propagation analysis.
Label mixing means that representations from examples with different labels overlap under a chosen projection or summary.
Such overlap can follow from poor feature learning, an unsuitable projection, or a task readout that does not use the available information.
It does not by itself show that repeated graph propagation made node representations converge.

A useful diagnostic panel can combine within-graph pairwise distances, cosine similarities, projection distance to a predicted invariant subspace, and task-conditioned probes.
Layerwise measurements show whether contraction develops before the task fails.
Comparisons across node roles distinguish general graph-wide contraction from loss of a particular source-to-target distinction.
The diagnostic should report the graph operator, normalisation, layer, node set, and scale convention because these choices change the trace being measured.

Attention does not remove the need for this test.
Attention-based graph systems can also contract under explicit assumptions, even though their aggregation weights are learned rather than fixed \cite{wu2023demystifying}.
The relevant question remains whether the learned propagation erases a task-relevant distinction.

\subsection{Model and Intervention Families}

The first response is to make the required depth usable without allowing repeated propagation to dominate the representation.
PairNorm maintains the spread of node representations and directly targets a geometric collapse statistic \cite{zhao2020pairnorm}.
Batch normalisation changes representation geometry and scale, while analyses of linearised GNNs show that normalisation can alter the limiting subspace rather than merely delaying convergence \cite{ioffe2015batch,scholkemper2025residual}.
These methods preserve variation in a general sense, but they do not select the source information that the target function requires.

Stochastic regularisation changes how often the same graph mixing is applied during training.
DropEdge randomly removes edges and can slow repeated smoothing while regularising a deep GNN \cite{rong2020dropedge}.
On a path-dependent task, the same removal can also delete an edge that carries useful information.
Its role must therefore be evaluated against both representation contraction and preservation of the required path.

Residual and dense connections give deep layers access to earlier representations.
DeepGCN-style residual and dense architectures preserve shallow features and provide shorter routes through the computation \cite{li2019deepgcns,he2015resnet}.
GCNII repeatedly introduces the initial node features and uses identity mapping to constrain how far the representation can move from its input \cite{chen2020simple}.
In linearised analyses, initial residuals constrain the attainable subspace, while batch normalisation can prevent one-dimensional collapse under the stated assumptions \cite{scholkemper2025residual}.

Layer-selection methods avoid forcing the task readout to use only the deepest representation.
Jumping Knowledge Networks combine intermediate layers and let the readout select among receptive-field sizes \cite{xu2018jknet}.
DAGNN likewise aggregates information from several depths \cite{liu2020dagnn}.
These methods can retain a useful shallow representation, but their selected layer may precede the arrival of the required source information.

Propagation-weighting methods replace a rigid stack with a mixture of propagation scales.
APPNP applies personalised PageRank-style propagation after feature transformation, while GPRGNN learns the weights assigned to different propagation steps \cite{gasteiger2019appnp,chien2021adaptive}.
Implicit and infinite-depth models such as IGNN, EIGNN, and graph implicit nonlinear diffusion use fixed-point or diffusion-like formulations instead of an explicit finite stack \cite{gu2020implicit,liu2021eignn,chen2022optimization}.
These systems can improve stability or access to several scales, but the learned mixture still needs to retain the task-relevant path information.

\subsection{What Successful Intervention Would Show}

A successful anti-smoothing intervention would preserve the relevant representation distinction as depth or dependency length increases and would improve the corresponding task outcome.
When the intervention directly controls the measured contraction, this joint change shows that representation geometry matters in the tested setting.
A stronger causal test would match capacity and optimisation, then show that the performance gain tracks restoration of the task-relevant distinction rather than a generic increase in variance.

Success from layer selection would establish that access to a suitable propagation depth is useful.
Success from residual or normalisation methods would establish that their combined change to representation flow and training makes the dependency easier to learn.
Neither result alone identifies over-smoothing because these interventions also alter optimisation, effective depth, and the hypothesis class available to the readout.

\subsection{What Remains Unproven}

Preserved variance does not prove that the representation retains source identity, path state, or any other feature needed by the label.
A model can avoid convergence while discarding the relevant information in a different direction.
Conversely, strong contraction need not be harmful when the target is constant on the contracted subspace.

An accuracy gain from an anti-smoothing method does not show that over-smoothing caused the baseline failure.
The method may improve gradient flow, regularise the model, expose shallower features, or change the balance of propagation scales.
A diagnosis requires the predicted contraction trace and a task-aligned link from that trace to the failure.

\section{Over-Squashing}

\subsection{Definition}

Over-squashing occurs when relevant information is structurally reachable but has too little usable influence on the target.
Whether the available computational channels are adequate depends on both the model and the task.
This definition excludes strict under-reaching, where the source lies outside the receptive field and its influence is exactly zero.
It also avoids reducing the mechanism to one graph shape or one growth law.

The literature supports four related accounts.
A topological account locates narrow cuts, unfavourable curvature, or other graph properties that restrict communication between relevant regions \cite{topping2022curvature}.
A path-proliferation account considers the pressure created when many relevant paths or message-passing walks feed a bounded target representation \cite{alon2021bottleneck}.
A finite-width account asks whether bounded representations and message functions can retain or select all relevant information as width, depth, and topology vary \cite{digionvanni2023oversquashing}.
A dynamical account measures decay in source-to-target sensitivity or task-relative information mixing under the learned propagation \cite{digionvanni2024power,bamberger2025measuring}.

These accounts overlap but are not equivalent.
Exponential growth in the number of paths or nodes is one setting that can create pressure, not the definition of over-squashing.
A graph with few paths can still show sensitivity decay along a long chain.
A graph with many paths need not fail when the task needs only a compressible summary.
The mechanism depends on which source information matters and on how the model carries it.

\subsection{Predicted Internal Trace}

A topological diagnosis predicts that task-relevant communication crosses a structural bottleneck whose severity increases in the failing regime.
Cuts, curvature, spectral connectivity, and path or walk counts can describe this structure.
These quantities are proxies because they do not include learned weights, source relevance, direction, or the target function.

A sensitivity diagnosis predicts that a perturbation at a relevant source node has little effect on the target representation or output even though the source is reachable.
Task-conditioned input-output Jacobians can measure this effect, and distance-weighted derivatives can summarise how model influence changes with range \cite{bamberger2025measuring}.
Perturbation tests provide a finite-difference counterpart that can be easier to interpret when the target function specifies which input change should matter.

A finite-width diagnosis predicts an interaction between communication demand and representation capacity.
With topology, depth, and training held fixed, increasing width should improve source retention or task performance when capacity is the limiting factor.
Distance-versus-branching controls can separate a long narrow chain from a shorter computation that aggregates many relevant sources.
Without such controls, a gain from additional width may reflect a generally larger model rather than relief of over-squashing.

Task-relative mixing predicts loss of the distinctions needed for the target even when a topology-only proxy appears favourable.
The diagnostic should identify the relevant source variables and test whether their influence remains separately usable at the target.
This criterion prevents irrelevant information in a large receptive field from being counted as communication success.

\subsection{Model and Intervention Families}

Diffusion and higher-order propagation broaden the communication operator so that distant nodes can interact without one learned layer per original edge.
DIGL constructs a diffused graph using operators such as personalised PageRank, while MixHop aggregates several powers of the adjacency matrix in separate channels \cite{gasteiger2019diffusion,abu2019mixhop}.
FA-GCN similarly changes connectivity under the bottleneck account of local message passing \cite{alon2021bottleneck,zhou2025glora}.
These methods can reduce the depth and sensitivity cost of distance, but they may expose endpoints without preserving which intermediate nodes make the dependency valid.

Shortest-path methods add explicit distance or path relations.
SP-GCN variants pass messages through shortest-path-based relations, which can make distant information available in fewer layers \cite{abboud2022shortest}.
This intervention changes the hypothesis class and the effective dependency relative to one-hop propagation on the original graph.
Its result therefore measures learning with shortest-path access rather than only improved transport along the original path.

Rewiring methods alter the communication graph before or during propagation.
Geom-GCN defines neighbourhoods from geometric structure, while DRew changes connectivity dynamically and controls when messages arrive across layers \cite{pei2020geomgcn,gutteridge2023drew}.
SDRF uses graph curvature to locate and relieve structural bottlenecks, and FOSR changes spectral connectivity through first-order rewiring \cite{topping2022curvature,karhadkar2023fosr}.
These methods directly target topology, but every new route may also shorten the effective dependency or create a shortcut to the label.

Global attention removes the one-hop visibility limit by allowing broad or all-pairs interactions.
GraphTransformer adapts transformer attention to graphs, SAN supplies spectral structure to attention, and GraphGPS combines local message passing with global attention \cite{vaswani2017attention,dwivedi2020graphtransformer,kreuzer2021san,rampasek2022gps}.
This family can bypass a topological bottleneck, but visibility alone does not tell the model which node, path, or relation matters.

Positional and structural encodings address this missing graph information.
Laplacian encodings provide spectral coordinates, random-walk encodings describe local-to-medium-range structure, and shortest-path or edge encodings expose selected relations.
GraphGPS uses Laplacian positional and random-walk structural encodings alongside local and global interaction channels \cite{rampasek2022gps}.
An encoding can support the intended relation, but it can also correlate with the label and become a shortcut.

\subsection{What Successful Intervention Would Show}

A successful communication intervention would increase task-relevant source-to-target influence and improve the task outcome under a matched comparison.
If a width increase restores the predicted sensitivity without changing topology or depth, the result supports a finite-capacity account in that regime.
If a targeted rewiring relieves a measured bottleneck while preserving the target dependency, the joint structural and behavioural change supports the corresponding topological account.

Success from global attention would show that a broader interaction channel helps the tested task.
Success from a structural encoding would show that the supplied graph relation helps the model select useful interactions.
These results establish practical value for the tested architecture, but they do not by themselves show that the baseline failed through over-squashing.

\subsection{What Remains Unproven}

Topology-only measures do not establish that the task requires communication across the measured bottleneck.
Many theoretical bounds also assume undirected graphs or particular derivative bounds, so their conclusions do not transfer automatically to directed learned systems.
Path counts, curvature, and spectral connectivity should therefore be reported with the graph direction, task, and model assumptions that make them relevant.

Rewiring can improve accuracy by changing the task presented to the model.
New edges may shorten the effective dependency, bypass required intermediate nodes, or expose an easier statistic.
Diffusion and higher-order propagation can likewise blur which original path carried the information.
A gain does not support the motivating mechanism unless the comparison preserves the dependency that the baseline failed to learn.

Global visibility does not prove task-relative information mixing or path-aware use.
Attention may connect the source and target while failing to represent the relevant intermediate structure.
An encoding may solve a structural identification problem without improving information transport, or it may leak a shortcut.
The experiment must still show that the intervention restores the specific usable influence predicted by the proposed account.

\section{Optimisation and Vanishing Gradients}

\subsection{Definition}

Optimisation failure is the broad case in which a model can express the required computation and has structural access to its inputs, but training does not find parameters that implement it.
Vanishing gradients are one specific optimisation mechanism.
They occur when repeated Jacobian products make the backward signal very small in the layers or parameters that must learn the computation \cite{bengio1994learning,glorot2010difficulty}.

For a path-dependent message-passing task, the loss is evaluated at the target output while early layers first process the source information.
The backward signal must cross the task readout and the intermediate message-passing updates before it can train those early transformations.
Greater depth can therefore weaken the signal even when the forward receptive field is large enough.

Optimisation and vanishing gradients must not be treated as synonyms.
Training can fail with gradients of ordinary magnitude because the objective has unhelpful directions, because gradients from different examples cancel, or because the inductive bias favours another function.
Conversely, a small aggregate gradient can occur near a useful optimum.
The diagnosis concerns whether task-aligned learning signals disappear where the required computation must be learned.

\subsection{Predicted Internal Trace}

A simple vanishing-gradient account predicts that gradient magnitudes decrease towards the early message-passing layers as depth or dependency length grows.
The trace should distinguish successful short settings from failing long settings under the same optimiser and measurement convention.
Layerwise norms over training are more informative than a single final value because vanishing may be transient or confined to the period when the dependency should be learned.

Parameter-gradient magnitude alone is not a measure of useful optimisation.
A nonzero weight gradient can be dominated by local features or regularisation terms.
Separately, gradients from different examples can cancel during aggregation.
Neither observation shows that source-conditioned information reaches every parameter required by the intended computation.
Joint analyses of graph and recurrent learning make this distinction explicit by relating particular contractive forward dynamics to particular backward-gradient effects \cite{arroyo2025vanishing}.

A task-aligned diagnostic should therefore pair gradient norms with the role of the parameter and the source-conditioned objective.
Possible tests include gradients of a target contrast under a controlled source perturbation and layerwise gradient alignment between examples that require the same path operation.
Intervention checks can then test whether restored early-layer updates improve use of the source information.
Training and validation curves can identify an optimisation gap, but they do not determine which backward mechanism produced it.

\subsection{Model and Intervention Families}

Residual and dense connections shorten backward routes and retain earlier representations.
DeepGCN-style systems use these connections to train many graph-convolution layers, while GCNII adds initial residuals and identity mapping \cite{li2019deepgcns,chen2020simple}.
Normalisation, suitable initialisation, and non-saturating activations can also stabilise scales across a deep computation \cite{ioffe2015batch,glorot2010difficulty}.
Improved gradient flow can make even simple GCN variants substantially stronger under some protocols \cite{jaiswal2022old}.

Gating gives the model learned control over retention and update strength.
GGNN uses recurrent-style gates for node updates, while GatedGCN applies gates to edge-level message passing \cite{li2016ggnn,bresson2017residual}.
Gradient-gating methods control propagation rates and backward signals more directly in deep graph learning \cite{rusch2023gradient}.
These interventions address trainability and information retention together rather than isolating one mechanism.

Layer selection and decoupled propagation can reduce the number of learned transformations through which gradients must pass.
Jumping Knowledge readouts, APPNP, GPRGNN, and DAGNN provide shorter or weighted routes from intermediate propagation states to the loss \cite{xu2018jknet,gasteiger2019appnp,chien2021adaptive,liu2020dagnn}.
Implicit formulations replace explicit unrolling with a fixed-point objective and require their own stability and optimisation analysis \cite{gu2020implicit,liu2021eignn,chen2022optimization}.

Adaptive optimisers, gradient clipping, and learning-rate schedules can improve numerical training without targeting a graph-specific mechanism.
Their inclusion in a protocol helps prevent avoidable instability.
Their success shows that the protocol matters, not that a particular forward failure has been diagnosed.

\subsection{What Successful Intervention Would Show}

A successful optimisation intervention would restore task-aligned signal in the affected layers and improve learning of the intended dependency under a matched protocol.
A residual or gating intervention is most informative when the experiment shows when early-layer gradients recover and how that recovery changes source-conditioned behaviour.
The result would support the claim that trainability limited the baseline in the tested regime.

An optimiser or initialisation change can establish protocol sensitivity.
A layer-selection method can establish that a shorter route from an intermediate state to the loss is useful.
These findings may guide model design even when they do not isolate vanishing gradients as the original cause.

\subsection{What Remains Unproven}

Nonzero gradients do not prove that optimisation is useful or task-aligned.
The gradients may point towards a function that fits local correlations, may fail to preserve source information, or may update parameters that are not decisive for the dependency.
Large norms can also indicate instability rather than productive learning.

A gain from residual connections, normalisation, or gating does not isolate a backward mechanism because each intervention also changes the forward representation dynamics.
Likewise, failure after applying common training safeguards does not rule out every optimisation problem.
The bounded conclusion concerns the measured gradient trace and the tested intervention, not optimisation in general.

\section{Interactions Among the Mechanisms}

The three mechanisms describe different parts of one learned computation, so they can occur together.
Over-smoothing and over-squashing concern the forward route from source information to the target representation.
Vanishing gradients concern the backward route from the loss to the parameters that shape that forward computation.
Contractive forward dynamics can therefore reduce both source sensitivity and the gradients that train early transformations under particular assumptions \cite{arroyo2025vanishing}.

Forward mechanisms can also reinforce each other without becoming identical.
Repeated mixing may contract differences between nodes while a finite-width representation loses the identity of one relevant source among many inputs.
A target representation can retain substantial overall variation yet lose the particular source-conditioned distinction needed for the label.
It can also become smooth even when the graph contains no severe topological bottleneck.

Interventions inherit the same overlap.
Residual and dense connections affect representation preservation and backward routes.
Normalisation affects scale, limiting subspaces, and optimisation.
Gating controls what enters a representation and how gradients pass through repeated updates.
Rewiring changes topology, effective depth, and the hypothesis class, while global attention changes both communication routes and structural selection demands.

These shared effects explain why intervention success cannot identify a mechanism by category label alone.
A residual model may improve by preserving shallow features, easing training, or allowing its readout to use a different effective depth.
Rewiring may instead relieve a bottleneck or simply shorten the dependency.
Gating may filter irrelevant messages, retain source state, or change gradient flow.
Mechanism identification requires a predicted trace and a comparison designed to separate these alternatives.

\section{From Candidate Mechanisms to Diagnostic Tests}

The diagnostic sequence begins with the task rather than the mechanism name.
The benchmark must establish which source information the target function requires and how dependency length $d$ is controlled.
An accuracy curve can then locate a failing regime, after which strict under-reaching and each candidate mechanism receive separate tests.
The final interpretation should state what was measured, which prediction it tests, and which alternatives remain open.

Table~\ref{tab:failure-diagnostic-synthesis} summarises this taxonomy.
Each row distinguishes the predicted trace from an intervention and limits the conclusion that intervention success can support.

\begin{table}[t]
\centering
\scriptsize
\setlength{\tabcolsep}{2.5pt}
\begin{tabular}{L{0.14\textwidth}L{0.20\textwidth}L{0.18\textwidth}L{0.20\textwidth}L{0.18\textwidth}}
\hline
\textbf{Mechanism or boundary} & \textbf{Predicted trace} & \textbf{Intervention family} & \textbf{What success would establish} & \textbf{What remains unproven} \\
\hline
Strict under-reaching & Source influence is exactly zero because the receptive field excludes it & More local layers, higher-order propagation, or global interaction & The architecture can access the source & The learned prediction uses the required dependency \\
Over-smoothing & Task-relevant representations contract towards a predicted invariant subspace or lose required separation & Normalisation, stochastic regularisation, residual or dense connections, or layer selection & Preserving the measured distinction helps in the tested regime & Contraction caused the baseline error or the preserved variation is task-relevant \\
Topological or capacity over-squashing & Relevant communication crosses a worsening bottleneck or improves with matched width & Rewiring, diffusion, shortest-path methods, or wider messages & Relieving the measured communication limit helps while controls are held fixed & The original dependency was preserved or the proxy diagnoses learned information flow \\
Sensitivity or task-relative over-squashing & Reachable source perturbations have weak target influence or lose task-relevant identity & Global attention, structural encodings, task-aware routing, or gating & The intervention restores usable source-to-target influence & Topology alone caused the loss or visibility implies path-aware use \\
Optimisation and vanishing gradients & Task-aligned backward signals decay in the layers or parameters that must learn the dependency & Initialisation, normalisation, residual or dense connections, gating, or optimiser changes & Restored training signal helps learn the tested dependency & Nonzero or large gradients are useful, correctly directed, or sufficient \\
\hline
\end{tabular}
\caption{Diagnostic taxonomy for long-range learning failure.
Rows link each mechanism or boundary to its predicted trace, applicable intervention families, and the conclusions that a successful intervention would and would not support.}
\label{tab:failure-diagnostic-synthesis}
\end{table}

The taxonomy bounds what Chapter 6 must test.
An over-smoothing test must specify the representations, layers, nodes, comparison set, and contraction measure.
An over-squashing test must name the relevant source information and distinguish topology, path proliferation, finite-width capacity, sensitivity decay, and task-relative mixing.
An optimisation test must identify the parameters and training period measured, then distinguish nonzero gradients from useful source-conditioned updates.

Intervention studies need comparable bounds.
A rewiring comparison must state whether the intervention changes the effective dependency.
A global-attention or encoding comparison must test structural use rather than visibility alone.
A residual, normalisation, or gating comparison must acknowledge its simultaneous effects on forward representations and backward optimisation.

The present chapter supplies only the predictions and evidential limits.
Low accuracy remains a symptom, and a mechanism becomes a supported explanation only when its task-aligned trace is present under the assumptions of the diagnostic.
Chapter~5 next develops the benchmark controls required to interpret these diagnostics by ruling out alternative interpretations of observed performance and internal traces.
Chapter~6 then applies both those controls and the diagnostic tests to the controlled setting provided by GLoRa and reports the resulting evidence \cite{zhou2025glora}.
